% questions.tex —— 1977 年·湖南卷 · 题目内容片段（无答案）
% 由 tools/generate_projects.py 依来源总账自 1977/1977_数学真题_湖南_试题.md 生成，勿手改（改源 MD 后重跑）


\gkpart{湖南省高等学校招生数学试题答案和评分标准}

\begin{gktip}
\textbf{来源说明}：本卷依据固定源 PDF《湖南省高等学校招生数学试题答案和评分标准》第 1—9 页（书籍页码 49—57）转录。\textbf{固定源是答案和评分标准册：原卷只载卷首说明、逐题解答与评分标准，未印题面}，题面无法逐题转录；以下题面由编者从源解答反推整理，每题题后的「源注」写明反推依据与未能确认之处。卷首「说明」称：各试题的解答方法和步骤是多种多样的，考生只要解答正确，可评满分；部分解答合理，酌情给分；解答和评分标准各地可参照研究实行。分值按源卷标注：第 1 题理科每小题 5 分、6 小题共 30 分，文科每小题 6 分、6 小题共 36 分；第 2—5 题理科各 10 分，文科依次为 12、12、13、12 分；第 6、7 题文理均 15 分。理科合计 100 分；文科按源卷标注合计 115 分，源卷未说明缘由。末尾「参考题」2 道不记分（但要作出评语），不计入题量。
\textbf{题号说明}：原卷正式题用中文序号「一、」至「七、」编排，本库改作阿拉伯数字第 1—7 题；第 1 题的 6 个小题原印作「1．」至「6．」，本库按小问编号规范重排为（1）—（6），未增删可确认的题面内容。
\end{gktip}

\par\medskip

\begin{question}
\gkti{1}

\Needspace{5\baselineskip}
（1） 解方程：

\nopagebreak[4]
\[
x^{-\frac{3}{2}}=\frac{27}{8}.
\]

（2） 计算：$2\lg6+2\sqrt{(\lg6-1)^2}$。

（3） 求 $x$ 的取值范围，$x$ 应使 $3-x\ge0$ 且 $x-2>0$。

\Needspace{5\baselineskip}
（4） 解不等式：

\nopagebreak[4]
\[
(x+4)(x-7)<0.
\]

\Needspace{5\baselineskip}
（5） 证明：

\nopagebreak[4]
\[
\cos15^\circ+\sin15^\circ=\frac{\sqrt6}{2}.
\]

（6） 求中心在原点、对称轴为坐标轴，且经过点 $(5,0)$ 与 $\left(4,\dfrac{12}{5}\right)$ 的椭圆方程。

\begin{gktip}
\textbf{源注}：源 PDF 未印题面。第（1）小题由源解答两种解法的首式与卷首注「由 $x^{-\frac32}=\dfrac{27}{8}$ 直接写出 $x=\dfrac49$ 者给满分」确认；第（5）小题由三种证法的共同起点确认；第（6）小题由源解答「设椭园的方程是 $\dfrac{x^2}{a^2}+\dfrac{y^2}{b^2}=1$，则 $\dfrac{5^2}{a^2}+\dfrac{0^2}{b^2}=1$，$\dfrac{4^2}{a^2}+\dfrac{\left(\frac{12}{5}\right)^2}{b^2}=1$」的代点过程确认（源卷以「椭园」指「椭圆」，照录说明）。第（2）（3）（4）小题源解答各只有一步：首小题写作「原式＝$2\lg6+2\sqrt{(\lg6-1)^2}$」，次小题写作「$x$ 所取的值应使 $3-x\ge0$，$x-2>0$」，第三小题写作「可化为 $(x+4)(x-7)<0$」——被计算的原始式子、原函数或原不等式的完整形式\textbf{无法从源 PDF 确认}，上述式子取自源解答的等号右端，题目按此照录。
\end{gktip}

\end{question}
\begin{question}
\gkti{2}某车间有 $86$ 名工人，加工甲、乙、丙三种部件。在同一时间内，每人可加工甲种部件 $15$ 件，或乙种部件 $12$ 件，或丙种部件 $9$ 件；一套部件由甲种 $3$ 件、乙种 $2$ 件、丙种 $1$ 件配成。应安排加工甲、乙、丙三种部件各多少人，才能使三种部件按要求配套？

\begin{gktip}
\textbf{源注}：源 PDF 未印题面，本题面由源解答反推：源解答设「安排加工甲种部件的 $x$ 人，乙种部件 $y$ 人，丙种部件 $z$ 人，能使加工的三种部件按要求配套」，列出总人数为 $86$、甲种总件数 $15x$ 等于丙种总件数 $9z$ 的 $3$ 倍、乙种总件数 $12y$ 等于丙种总件数的 $2$ 倍三个等量关系，故总人数、三种效率 $15$、$12$、$9$ 与配套比 $3:2:1$ 可以确认；「车间」「同一时间内」等措辞为整理时的表述，原句无法确认。
\end{gktip}

\end{question}
\begin{question}
\gkti{3}已知关于 $x$ 的方程

\nopagebreak[4]
\[
ax^2-(2a+1)x+(a+1)=0
\]

（$a\ne0$）。$a$ 取何值时：

（1） 方程有一个根为零；

（2） 方程的两根互为相反数；

（3） 方程的两根之积等于 $2$。

\begin{gktip}
\textbf{源注}：源 PDF 未印题面，本题面由源解答反推：源解答由求根公式得两根 $x_1=\dfrac{a+1}{a}$、$x_2=1$，再按「若 $x_1=0$」「若 $x_1+x_2=0$」「若 $x_1x_2=2$」三种情形定 $a$，方程与三种情形均可确认；「关于 $x$ 的方程」等措辞为整理时的表述。情形（3）原卷也可能问「有一根为 $2$」——因两根中恒有一根为 $1$，两种问法等价。
\end{gktip}

\end{question}
\begin{question}
\gkti{4}有三个小锡球，表面积分别为 $36\pi$、$64\pi$、$100\pi$ 平方厘米。把这三个小锡球熔化后铸成一个大锡球，求大锡球的半径。

\begin{gktip}
\textbf{源注}：源 PDF 未印题面，本题面由源解答反推：源解答设三个小锡球半径依次为 $r_1$、$r_2$、$r_3$ 厘米、大锡球半径为 $R$ 厘米，从 $4\pi r_1^2=36\pi$ 等三式读出三个表面积，从「三个球体积之和等于大球体积」的等量关系读出熔铸过程；「平方厘米」「大锡球」由「厘米」的单位标注确认。
\end{gktip}

\end{question}
\begin{question}
\gkti{5}四边形 $ABCD$ 的四个内角 $\angle A$、$\angle B$、$\angle C$、$\angle D$ 依次成等差数列，且 $\angle C=3\angle A$，求四个内角。

\begin{gktip}
\textbf{源注}：源 PDF 未印题面，本题面由源解答反推：源解答设 $\angle A=x-3d$、$\angle B=x-d$、$\angle C=x+d$、$\angle D=x+3d$（即四角成等差数列），并用「$\angle C=3\angle A$」列出第二方程，可确认题设只有这两个条件；「四边形 $ABCD$」为整理时的表述。
\end{gktip}

\end{question}
\begin{question}
\gkti{6}在 $\triangle ABC$ 中，已知 $\angle B=45^\circ$，$b=10$，$c=5\sqrt6$，解此三角形。

\begin{gktip}
\textbf{源注}：源 PDF 未印题面，本题面由源解答反推：源解答首步即正弦定理式 $\dfrac{10}{\sin45^\circ}=\dfrac{5\sqrt6}{\sin C}$，据此可知已知角为 $45^\circ$、其对边为 $10$、另一已知边 $5\sqrt6$ 的对角为 $C$；标准边角记号 $b$、$c$ 为整理时的表述。
\end{gktip}

\end{question}
\begin{question}
\gkti{7}圆内接四边形 $ABCD$ 中，$\angle A=60^\circ$，$\angle B=90^\circ$，$AB=2$，$CD=1$，求 $BC$ 和 $DA$。

\begin{gktip}
\textbf{源注}：源 PDF 未印题面，本题面由源解答反推：「圆内接」由源解答首步「$\angle A+\angle C=\angle B+\angle D=180^\circ$」确认；解法一延长 $BC$、$AD$ 交于 $E$ 后得 $AE=4$、$CE=2$，并用 $BE=\sqrt{4^2-2^2}$、$DE=\sqrt{2^2-1}$，据此读出 $AB=2$、$CD=1$；解法二所列两式 $x^2+2^2=y^2+1^2$、$x^2+1^2-2x\cos120^\circ=y^2+2^2-4y\cos60^\circ$ 中的 $2$、$1$ 与 $\angle C=120^\circ$、$\angle A=60^\circ$ 同此。源卷此题无附图。
\textbf{参考题}（源卷注明：不记分，但要作出评语）
\end{gktip}

\Needspace{5\baselineskip}
\textbf{参考题 1．} 用数学归纳法或其他方法证明：

\nopagebreak[4]
\[
1^4+2^4+3^4+\cdots+n^4=\frac{n(n+1)(2n+1)(3n^2+3n-1)}{30}.
\]

\textbf{参考题 2．} 计算下列积分：

（1）$\int_0^1 x\mathrm e^{-x}\,dx$；

（2）$\int_0^a x^2\sqrt{a^2-x^2}\,dx$。

\end{question}